\section{Statistical Methods}
\label{sec:statistical-methods}

\subsection{The Performance Score (z-score)}
Each participant's reported result \(x_i\) for a given analyte is converted to a standardised
performance score using the classical z-score formulation of ISO 13528 \S 9:

\begin{equation}
z_i \;=\; \frac{x_i - x_{pt}}{\sigma_{pt}}
\end{equation}

\noindent where \(x_{pt}\) is the assigned value and \(\sigma_{pt}\) is the standard deviation for
proficiency assessment, both defined in Section~\ref{sec:scheme-design}. The z-score expresses a
participant's deviation from the assigned value in units of the fitness-for-purpose target
standard deviation, allowing performance to be compared directly across analytes with very
different concentration ranges and units.

\subsection{Performance Classification}
Following the convention adopted throughout ISO 13528 and the great majority of accredited PT
schemes, results are classified into three performance bands:

\begin{center}
\renewcommand{\arraystretch}{1.4}
\begin{tabular}{@{}lll@{}}
\toprule
\textbf{Score range} & \textbf{Classification} & \textbf{Interpretation} \\
\midrule
\rowcolor{ptteal!10} \(|z| \leq 2\) & \perfsat & No action required. \\
\rowcolor{ptamber!14} \(2 < |z| < 3\) & \perfque & Warning signal; review recommended. \\
\rowcolor{ptred!10} \(|z| \geq 3\) & \perfuns & Action signal; investigation required. \\
\bottomrule
\end{tabular}
\end{center}

This banding is the fitness-for-purpose convention widely adopted for chemical-analysis PT
schemes following Thompson's analysis of achievable inter-laboratory precision at trace
concentrations \cite{thompson2000zscores}. A single Questionable score is treated as an informal early-warning signal. Two Questionable
scores for the same analyte in three consecutive rounds, or a single Unsatisfactory score, trigger
the formal corrective-action process described in Section~\ref{sec:recommendations}.

\subsection{Robustness of the Assigned Value}
Because \(x_{pt}\) is computed from the participant population itself (Section~\ref{sec:scheme-design}),
Algorithm A down-weights outlying results iteratively so that no single laboratory's result can
materially shift the assigned value against which it is itself judged. The robust standard
deviation produced by Algorithm A (\(s^{*} = 0.61\)~\(\mu\)g/L for Pb this round) was additionally
checked against the fixed fitness-for-purpose \(\sigma_{pt}\) and found to differ by less than
10\%, confirming that the fitness-for-purpose target remains realistic for the current state of
the art among participating laboratories.
